\documentclass[10pt]{article}
\usepackage[T1]{fontenc}
\usepackage[utf8]{inputenc}
\usepackage[french]{babel}
\usepackage{amsmath,amssymb,mathrsfs,mathpazo}
\usepackage{geometry}
\geometry{a4paper,margin=2cm}
\usepackage{graphicx,array,tabularx,tabularray,colortbl,multirow,multicol,booktabs,enumitem,pifont,ulem,fancybox,diagbox,cancel}
\usepackage{pstricks,pstricks-add,pst-plot,pst-tree,pst-node,pst-text,pst-func,pst-math,pst-bspline,pst-3dplot,pst-eucl,pst-solides3d}
\let\pscalculate\undefined
\usepackage{pst-bezier}
\usepackage{tikz,pgfplots,tkz-tab}
\usetikzlibrary{arrows,arrows.meta,patterns,calc,positioning,trees,decorations.markings,shapes.geometric,shapes.misc,shapes.arrows,intersections,decorations.pathmorphing,decorations.pathreplacing,backgrounds,shadows,babel}
\pgfplotsset{compat=1.18}
\usepgfplotslibrary{fillbetween}
\usepackage[np]{numprint}
\usepackage{listings,makecell,pas-tableur,twoopt}
\let\barre\undefined
\usepackage[pstricks]{tablvar}
\definecolor{darkgreen}{RGB}{0,100,0}
\providecommand{\red}{\color{red}}
\providecommand{\blue}{\color{blue}}
\providecommand{\green}{\color{green}}
\providecommand{\black}{\color{black}}
\providecommand{\white}{\color{white}}
\providecommand{\gray}{\color{gray}}
\providecommand{\cyan}{\color{cyan}}
\providecommand{\magenta}{\color{magenta}}
\providecommand{\yellow}{\color{yellow}}
\providecommand{\violet}{\color{violet}}
\providecommand{\orange}{\color{orange}}
\providecommand{\euro}{\texteuro}
\providecommand{\EUR}{\texteuro}
\providecommand{\ieme}{\textsuperscript{e}}
\providecommand{\degres}{\ensuremath{^\circ}}
\providecommand{\degree}{\ensuremath{^\circ}}
\providecommand{\vv}[1]{\overrightarrow{#1}}
\providecommand{\Coord}[1]{#1}
\providecommand{\decosix}{}
\providecommand{\starredbullet}{\ensuremath{\bullet}}
\pagestyle{empty}
\setlength{\parindent}{0pt}
\usepackage{adjustbox}\DeclareUnicodeCharacter{2236}{:}\DeclareUnicodeCharacter{2190}{\ensuremath{\leftarrow}}
\providecommand{\R}{\mathbb{R}}
\providecommand{\N}{\mathbb{N}}
\providecommand{\Z}{\mathbb{Z}}
\providecommand{\Q}{\mathbb{Q}}
\providecommand{\D}{\mathbb{D}}
\providecommand{\C}{\mathbb{C}}
\providecommand{\e}{{\mathrm{e}}}
\providecommand{\vect}[1]{\overrightarrow{#1}}
\providecommand{\vectt}[1]{\overrightarrow{\mathrm{#1}}}
\providecommand{\vvt}[1]{\overrightarrow{\mathrm{#1}}}
\providecommand{\np}[2][]{{#2~#1}}
\providecommand{\vv}[1]{{\overrightarrow{#1}}}
\renewcommand{\d}{{\mathrm{d}}}
\providecommand{\degres}{{^\circ}}
\providecommand{\degre}{{^\circ}}
\providecommand{\blue}{\color{blue}}
\providecommand{\red}{\color{red}}
\providecommand{\enskip}{\quad}
\providecommand{\llbracket}{\lbrack\!\lbrack}
\providecommand{\rrbracket}{\rbrack\!\rbrack}
\providecommand{\euro}{{\text{€}}}
\providecommand{\no}{{\text{n°}}}
\providecommand{\interval}[2]{{[#1;#2]}}
\providecommand{\intervFO}[2]{{[#1;#2[}}
\providecommand{\dotfill}{\ldots}

\providecommand{\ds}{\displaystyle}
\providecommand{\textding}[1]{#1}
\providecommand{\Oij}{\(\mathrm{O};\vec{\imath},\vec{\jmath}\)}
\providecommand{\Oijk}{\(\mathrm{O};\vec{\imath},\vec{\jmath},\vec{k}\)}
\providecommand{\Ouv}{\(\mathrm{O};\vec{u},\vec{v}\)}
\providecommand{\intervFF}[2]{\ensuremath{[#1;#2]}}
\renewcommand{\euro}{\ifmmode\text{\texteuro}\else\texteuro\fi}
\DeclareUnicodeCharacter{00B0}{\ifmmode^{\circ}\else\textdegree\fi}
\DeclareUnicodeCharacter{20AC}{\ifmmode\text{\texteuro}\else\texteuro\fi}
% Énoncés : Éducation nationale. Transcription et corrigés : APMEP, Denis Vergès et contributeurs. Découpage et classement : CoopMaths, Sébastien Lozano, Éric Elter et contributeurs. Adaptation de la présentation : Maths Explorer.
% CC BY-SA : https://coopmaths.fr/www/about/
% https://coopmaths.fr/alea/static/bac/2024/tex/bac_2024_06_sujet2_metropole_devoile_4.tex
\begin{document}
\noindent\textit{APMEP / CoopMaths. CC BY-SA. Présentation adaptée par Maths Explorer.}\par\medskip
\medskip

Soit $f$ une fonction définie et deux fois dérivable sur $\R$. On note $f'$ sa fonction dérivée et $f''$ sa dérivée seconde.

\begin{list}{\textbullet}{Dans le repère orthonormé ci-dessous ont été représentés :}
\item la courbe représentative $\mathcal{C}_f$ de la fonction $f$ ;
\item la tangente $T$ à $\mathcal{C}_f$ en son point N(0~;~2);
\item le point M$(-2~;~0)$ appartenant à $\mathcal{C}_f$ et P(2~;~0) appartenant à la tangente $T$.
\end{list}

On précise que la fonction $f$ est strictement positive sur l'intervalle $[0~;~+\infty[$ et qu'elle est strictement croissante sur l'intervalle $]- \infty~;~-1]$.

\begin{center}
\psset{unit=1cm,arrowsize=2pt 3}
\begin{pspicture*}(-2.5,-1.5)(5.5,3.5)
\psgrid[subgriddiv=1,  gridlabels=0, gridcolor=lightgray] 
\psaxes[linewidth=1.25pt,labelFontSize=\scriptstyle]{->}(0,0)(-2.5,-1.5)(5.5,3.5)
\psplot[plotpoints=3000,linewidth=1.25pt,linecolor=red]{-2.5}{5.5}{x 2 add 2.71828 x exp div}
\psplotTangent{0}{4.5}{x 2 add 2.71828 x exp div}
\psdots(-2,0)(0,2)(2,0)
\uput[ul](-2,0){M} \uput[ur](0,2){N} \uput[ur](2,0){P}\uput[d](3,-1){$T$}
\end{pspicture*}
\end{center}


\begin{center}\textbf{Partie A : étude graphique}\end{center}

On répondra aux questions suivantes en utilisant le graphique.

\begin{enumerate}
\item 
	\begin{enumerate}
		\item Donner $f(0)$.
		\item Déterminer $f'(0)$.
	\end{enumerate}
\item Résoudre l'équation $f(x) = 0$.
\item La fonction $f$ est-elle convexe sur $\R$ ? Justifier.
\item Parmi les courbes suivantes, indiquer laquelle peut représenter une primitive de la fonction $f$ sur $\R$. Justifier.

\begin{center}
\begin{tabularx}{\linewidth}{|*{3}{>{\centering \arraybackslash}X|}}\hline
Courbe 1&Courbe 2&Courbe 3\\ \hline
\psset{unit=0.4cm,arrowsize=2pt 3}
\begin{pspicture*}(-2.5,-8)(8,3)
\psaxes[linewidth=1.25pt,labelFontSize=\scriptstyle]{->}(0,0)(-2.5,-8)(8,3)
\psplot[plotpoints=3000,linewidth=1.25pt,linecolor=red]{-2.5}{5.5}{x 1 add 2.71828 x exp div neg}
\end{pspicture*}&
\psset{unit=0.4cm,arrowsize=2pt 3}
\begin{pspicture*}(-4,-8)(6,3)
\psaxes[linewidth=1.25pt,labelFontSize=\scriptstyle]{->}(0,0)(-4,-8)(6,3)
\psplot[plotpoints=3000,linewidth=1.25pt,linecolor=red]{-4}{5.5}{x 3 add 2.71828 x exp div neg}
\end{pspicture*}
&\psset{unit=0.4cm,arrowsize=2pt 3}
\begin{pspicture*}(-2.5,-8)(8,3)
\psaxes[linewidth=1.25pt,labelFontSize=\scriptstyle]{->}(0,0)(-2.5,-8)(8,3)
\psplot[plotpoints=3000,linewidth=1.25pt,linecolor=red]{-2.5}{5.5}{x 1 sub 2.71828 x exp div neg}
\end{pspicture*}\\ \hline
\end{tabularx}
\end{center}
\end{enumerate}

\medskip

\begin{center}\textbf{Partie B :  recherche d'une expression algébrique}\end{center}

On admet que la fonction $f$ est de la forme 
\[f(x)= (ax + b)\e^{\lambda x},\]
où $a, b$ et $\lambda$ sont des constantes réelles.

Pour répondre aux questions suivantes, on utilisera les résultats de la partie A.

\medskip

\begin{enumerate}
\item Justifier que $b = 2$.
\item Justifier que $-2a + b = 0$ puis en déduire la valeur de $a$.
\item Déterminer une expression algébrique de $f$. Justifier.
\end{enumerate}

\medskip

\begin{center}\textbf{Partie C : étude algébrique}\end{center}

On admet que la fonction $f$ est définie sur $\R$ par 
\[f(x) = (x + 2)\e^{-x}.\]

\smallskip

\begin{enumerate}
\item Déterminer la limite de $f$ en $-\infty$.
\item On admet que $f'(x) =(- x - 1)\e^{-x}$. Dresser le tableau de variations complet de $f$.

Justifier.
\item 
	\begin{enumerate}
		\item Étudier la convexité de $f$.
		\item Préciser les coordonnées des éventuels points d'inflexion de la courbe $\mathcal{C}_f$.
	\end{enumerate}
\item Pour tout nombre réel $t \geqslant 0$, on pose:
\[I(t) = \displaystyle\int_{-2}^t f(x)\:\text{d}x.\]

	\begin{enumerate}
		\item En utilisant une intégration par parties, montrer que:
\[I(t) = (- t - 3)\e^{-t} + \e^2.\]

		\item En déduire un exemple de surface non limitée dont l'aire est finie.
	\end{enumerate}
\end{enumerate}

\end{document}
